\documentclass[10pt]{article}
\usepackage[a4paper,margin=1.6cm]{geometry}
\usepackage{booktabs}
\usepackage{longtable}
\usepackage{amsmath}
\usepackage[T1]{fontenc}
\usepackage{mathptmx}
\usepackage{url}

\renewcommand{\thetable}{S\arabic{table}}
\setlength{\LTcapwidth}{\textwidth}

\title{Supplementary Material for\\[2pt]
``Beyond Leaderboards: Deployment-Context-Aware Evaluation of\\ Small and Large Language Models''}
\author{N. S. Stephane, G. Manzo, and P. Delgado}
\date{}

\begin{document}
\maketitle

\noindent
This document reports the uncertainty of the system-level measurements and the
sensitivity analyses referred to in the article. All values are computed from the
per-run result files of the benchmarking framework, available at
\url{https://github.com/pameladelgado/TB-2026-SLM-benchmarking}, with the script
\texttt{compute\_supplementary.py}; before computing intervals, the script
reproduces every mean reported in the article from these files (maximum relative
deviation 0.2\%). All measurements describe one sequential request stream.
Output length is counted as whitespace-separated words of the generated text, as
implemented by the framework.

For some configurations, measurements were collected in more than one session
(for example, after a re-run). As in the article, means pool the sessions of a
configuration with equal weight; intervals for timing metrics are obtained with
a two-stage cluster bootstrap (runs, then requests within runs; 2000 replicates)
stratified by session.

\begingroup\scriptsize\setlength{\tabcolsep}{3pt}
\begin{longtable}{@{}l l l c r r r r@{}}
\caption{Mean time to first token (TTFT), end-to-end latency, and per-request throughput with 95\% intervals (two-stage cluster bootstrap over runs and requests, stratified by measurement session, sessions weighted equally). Throughput is generated words per second of the decoding interval. S/R: number of measurement sessions and measured runs pooled.}\label{tab:s1}\\
\toprule
Context & Configuration & Workload & S/R & TTFT (ms) & Latency (ms) & Throughput (words/s) & Words \\
\midrule\endfirsthead
\toprule
Context & Configuration & Workload & S/R & TTFT (ms) & Latency (ms) & Throughput (words/s) & Words \\
\midrule\endhead
\bottomrule\endfoot
Edge & Gemma 3 1B (Q8) & MMLU & 1/10 & 1378 [1310, 1459] & 1566 [1494, 1637] & 5.41 [5.35, 5.47] & 1.0 \\
Edge & Gemma 3 1B (Q8) & GSM8K & 2/19 & 757 [747, 768] & 23524 [23322, 23711] & 6.11 [6.08, 6.15] & 139.4 \\
Edge & Gemma 3 1B (Q8) & HellaSwag & 1/10 & 1865 [1818, 1908] & 2056 [2012, 2097] & 5.28 [5.23, 5.33] & 1.0 \\
Edge & Gemma 3 1B (Q8) & Dolly & 2/10 & 1178 [1023, 1336] & 17449 [16454, 18414] & 6.51 [6.38, 6.62] & 108.1 \\
Edge & Llama 3.2 1B (Q4) & MMLU & 1/10 & 1693 [1589, 1799] & 1776 [1677, 1884] & 12.70 [12.54, 12.86] & 1.0 \\
Edge & Llama 3.2 1B (Q4) & GSM8K & 2/19 & 850 [838, 862] & 16013 [15828, 16188] & 9.49 [9.46, 9.53] & 144.1 \\
Edge & Llama 3.2 1B (Q4) & HellaSwag & 1/10 & 2347 [2295, 2400] & 2432 [2379, 2487] & 12.36 [12.21, 12.50] & 1.0 \\
Edge & Llama 3.2 1B (Q4) & Dolly & 2/10 & 1440 [1244, 1669] & 12098 [11234, 12980] & 9.29 [9.13, 9.44] & 100.4 \\
Edge & Phi-4 Mini 3.8B (Q4) & MMLU & 1/10 & 5798 [5454, 6168] & 6552 [6201, 6911] & 3.55 [3.49, 3.61] & 2.6 \\
Edge & Phi-4 Mini 3.8B (Q4) & GSM8K & 2/10 & 2873 [2778, 2975] & 42040 [40554, 43612] & 3.11 [3.07, 3.14] & 120.3 \\
Edge & Phi-4 Mini 3.8B (Q4) & HellaSwag & 1/10 & 8085 [7906, 8278] & 8862 [8679, 9047] & 3.63 [3.56, 3.70] & 3.0 \\
Edge & Phi-4 Mini 3.8B (Q4) & Dolly & 2/10 & 4940 [4234, 5722] & 38537 [36215, 40878] & 3.28 [3.23, 3.33] & 113.2 \\
Edge & SmolLM3 3B (Q4) & MMLU & 1/10 & 5051 [4765, 5364] & 5666 [5363, 5980] & 4.64 [4.52, 4.74] & 2.8 \\
Edge & SmolLM3 3B (Q4) & GSM8K & 2/10 & 2598 [2524, 2679] & 39711 [38821, 40551] & 4.03 [3.99, 4.08] & 149.1 \\
Edge & SmolLM3 3B (Q4) & HellaSwag & 1/10 & 6814 [6657, 6964] & 7445 [7284, 7599] & 4.88 [4.80, 4.96] & 3.1 \\
Edge & SmolLM3 3B (Q4) & Dolly & 2/10 & 4283 [3672, 4929] & 36428 [34706, 38175] & 4.06 [4.00, 4.12] & 133.5 \\
\midrule
Local CPU & Gemma 3 1B (Q8) & MMLU & 2/11 & 1048 [963, 1146] & 1128 [1043, 1222] & 13.25 [12.87, 13.65] & 1.0 \\
Local CPU & Gemma 3 1B (Q8) & GSM8K & 2/11 & 693 [662, 727] & 10607 [10322, 10886] & 14.41 [14.13, 14.71] & 139.5 \\
Local CPU & Gemma 3 1B (Q8) & HellaSwag & 2/11 & 1471 [1397, 1546] & 1552 [1479, 1631] & 12.81 [12.51, 13.11] & 1.0 \\
Local CPU & Gemma 3 1B (Q8) & Dolly & 2/11 & 1127 [971, 1282] & 8965 [8380, 9513] & 13.73 [13.22, 14.31] & 108.0 \\
Local CPU & Llama 3.2 1B (Q8) & MMLU & 1/10 & 1764 [1659, 1874] & 1821 [1718, 1937] & 18.14 [17.98, 18.29] & 1.0 \\
Local CPU & Llama 3.2 1B (Q8) & GSM8K & 1/10 & 934 [909, 960] & 3221 [3016, 3438] & 12.47 [12.23, 12.71] & 28.0 \\
Local CPU & Llama 3.2 1B (Q8) & HellaSwag & 1/10 & 2408 [2346, 2474] & 2464 [2400, 2529] & 18.02 [17.86, 18.16] & 1.0 \\
Local CPU & Llama 3.2 1B (Q8) & Dolly & 1/10 & 1694 [1461, 1981] & 8858 [8303, 9441] & 13.05 [12.80, 13.33] & 94.4 \\
Local CPU & Ministral 3 3B (Q4) & MMLU & 2/20 & 5581 [5351, 5822] & 5702 [5464, 5957] & 8.35 [8.29, 8.42] & 1.0 \\
Local CPU & Ministral 3 3B (Q4) & GSM8K & 2/11 & 2472 [2386, 2566] & 23060 [22418, 23729] & 5.76 [5.62, 5.90] & 116.6 \\
Local CPU & Ministral 3 3B (Q4) & HellaSwag & 2/20 & 7558 [7430, 7687] & 7681 [7553, 7813] & 8.21 [8.15, 8.27] & 1.0 \\
Local CPU & Ministral 3 3B (Q4) & Dolly & 2/11 & 3606 [3150, 4148] & 22395 [21514, 23322] & 6.43 [6.28, 6.58] & 118.1 \\
Local CPU & Phi-4 Mini 3.8B (Q4) & MMLU & 3/30 & 4079 [3933, 4237] & 4396 [4243, 4552] & 8.68 [8.55, 8.79] & 2.5 \\
Local CPU & Phi-4 Mini 3.8B (Q4) & GSM8K & 2/20 & 2412 [2366, 2456] & 18629 [18231, 19046] & 7.44 [7.31, 7.56] & 117.2 \\
Local CPU & Phi-4 Mini 3.8B (Q4) & HellaSwag & 2/20 & 6029 [5921, 6141] & 6391 [6275, 6500] & 8.28 [8.15, 8.40] & 3.0 \\
Local CPU & Phi-4 Mini 3.8B (Q4) & Dolly & 2/11 & 3613 [3168, 4098] & 18496 [17486, 19513] & 7.78 [7.60, 7.96] & 114.9 \\
\midrule
GPU server & Apertus 70B (bf16) & MMLU & 1/10 & 216 [213, 219] & 343 [340, 347] & 8.01 [7.98, 8.05] & 1.0 \\
GPU server & Apertus 70B (bf16) & GSM8K & 1/10 & 198 [198, 199] & 10969 [10801, 11139] & 12.03 [11.96, 12.11] & 129.4 \\
GPU server & Apertus 70B (bf16) & HellaSwag & 1/10 & 224 [222, 226] & 349 [347, 351] & 8.02 [7.97, 8.08] & 1.0 \\
GPU server & Apertus 70B (bf16) & Dolly & 1/10 & 213 [209, 219] & 10525 [10158, 10888] & 12.90 [12.70, 13.10] & 136.9 \\
GPU server & Granite 4 H-Tiny (bf16) & MMLU & 2/20 & 178 [177, 179] & 296 [295, 297] & 9.06 [8.97, 9.15] & 1.1 \\
GPU server & Granite 4 H-Tiny (bf16) & GSM8K & 2/20 & 185 [183, 186] & 7333 [7162, 7510] & 14.92 [14.65, 15.19] & 106.3 \\
GPU server & Granite 4 H-Tiny (bf16) & HellaSwag & 2/20 & 176 [176, 177] & 297 [295, 299] & 9.04 [8.95, 9.14] & 1.1 \\
GPU server & Granite 4 H-Tiny (bf16) & Dolly & 1/10 & 182 [180, 183] & 7324 [6978, 7670] & 15.67 [15.39, 15.99] & 117.0 \\
GPU server & Kimi Linear 48B-A3B (bf16) & MMLU & 1/10 & 183 [181, 184] & 315 [312, 318] & 14.58 [14.10, 15.11] & 2.0 \\
GPU server & Kimi Linear 48B-A3B (bf16) & GSM8K & 1/10 & 185 [182, 188] & 5875 [5617, 6130] & 15.41 [14.89, 16.01] & 89.4 \\
GPU server & Kimi Linear 48B-A3B (bf16) & HellaSwag & 1/10 & 183 [182, 184] & 308 [306, 311] & 13.46 [12.96, 13.98] & 1.8 \\
GPU server & Kimi Linear 48B-A3B (bf16) & Dolly & 1/10 & 182 [180, 184] & 5377 [5073, 5682] & 17.22 [16.60, 17.87] & 91.8 \\
GPU server & Llama 3.2 1B (bf16) & MMLU & 1/10 & 129 [128, 131] & 239 [237, 242] & 9.08 [8.99, 9.16] & 1.0 \\
GPU server & Llama 3.2 1B (bf16) & GSM8K & 1/10 & 129 [127, 130] & 1098 [1070, 1124] & 116.65 [114.30, 119.42] & 112.6 \\
GPU server & Llama 3.2 1B (bf16) & HellaSwag & 1/10 & 129 [128, 131] & 241 [239, 243] & 9.01 [8.92, 9.10] & 1.0 \\
GPU server & Llama 3.2 1B (bf16) & Dolly & 1/10 & 129 [128, 130] & 1120 [1072, 1166] & 116.89 [114.59, 119.07] & 120.8 \\
GPU server & Llama 3.3 70B (FP8) & MMLU & 2/20 & 181 [179, 182] & 296 [294, 297] & 8.70 [8.65, 8.75] & 1.0 \\
GPU server & Llama 3.3 70B (FP8) & GSM8K & 2/20 & 168 [167, 169] & 7058 [6989, 7126] & 20.95 [20.88, 21.01] & 143.1 \\
GPU server & Llama 3.3 70B (FP8) & HellaSwag & 2/20 & 190 [189, 192] & 306 [305, 307] & 8.66 [8.61, 8.70] & 1.0 \\
GPU server & Llama 3.3 70B (FP8) & Dolly & 1/10 & 182 [179, 186] & 6288 [5992, 6570] & 18.47 [18.22, 18.73] & 115.7 \\
GPU server & Ministral 3 14B & MMLU & 1/10 & 158 [157, 159] & 274 [272, 276] & 8.62 [8.55, 8.69] & 1.0 \\
GPU server & Ministral 3 14B & GSM8K & 1/10 & 158 [156, 159] & 11527 [11358, 11696] & 10.72 [10.55, 10.89] & 122.1 \\
GPU server & Ministral 3 14B & HellaSwag & 1/10 & 158 [157, 159] & 275 [273, 276] & 8.57 [8.51, 8.63] & 1.0 \\
GPU server & Ministral 3 14B & Dolly & 1/10 & 165 [162, 167] & 9962 [9416, 10557] & 11.83 [11.31, 12.30] & 115.7 \\
GPU server & Ministral 3 3B (AWQ) & MMLU & 1/10 & 130 [129, 132] & 242 [239, 244] & 9.02 [8.93, 9.10] & 1.0 \\
GPU server & Ministral 3 3B (AWQ) & GSM8K & 1/10 & 129 [128, 130] & 2917 [2821, 3005] & 44.79 [43.41, 46.42] & 125.1 \\
GPU server & Ministral 3 3B (AWQ) & HellaSwag & 1/10 & 132 [131, 133] & 243 [241, 245] & 9.01 [8.92, 9.10] & 1.0 \\
GPU server & Ministral 3 3B (AWQ) & Dolly & 1/10 & 132 [131, 134] & 2530 [2444, 2613] & 49.82 [48.45, 51.19] & 120.5 \\
GPU server & Ministral 3 8B (AWQ) & MMLU & 1/10 & 134 [133, 135] & 244 [243, 246] & 9.04 [8.99, 9.09] & 1.0 \\
GPU server & Ministral 3 8B (AWQ) & GSM8K & 1/10 & 131 [130, 132] & 3537 [3478, 3624] & 37.16 [36.32, 37.82] & 126.6 \\
GPU server & Ministral 3 8B (AWQ) & HellaSwag & 1/10 & 137 [136, 138] & 249 [248, 251] & 8.90 [8.84, 8.96] & 1.0 \\
GPU server & Ministral 3 8B (AWQ) & Dolly & 1/10 & 135 [134, 136] & 3202 [3126, 3276] & 38.92 [37.60, 40.23] & 120.3 \\
GPU server & Phi-4 14B (bf16) & MMLU & 1/10 & 147 [146, 149] & 261 [259, 264] & 11.41 [10.97, 11.87] & 1.3 \\
GPU server & Phi-4 14B (bf16) & GSM8K & 1/10 & 141 [140, 142] & 2861 [2821, 2906] & 49.96 [49.31, 50.43] & 135.1 \\
GPU server & Phi-4 14B (bf16) & HellaSwag & 1/10 & 150 [149, 151] & 263 [262, 266] & 9.85 [9.55, 10.19] & 1.1 \\
GPU server & Phi-4 14B (bf16) & Dolly & 1/10 & 144 [143, 146] & 2478 [2384, 2573] & 55.73 [55.03, 56.40] & 132.7 \\
GPU server & Phi-4 Mini 3.8B (GPTQ-4) & MMLU & 1/10 & 139 [139, 140] & 253 [252, 254] & 20.55 [19.93, 21.17] & 2.3 \\
GPU server & Phi-4 Mini 3.8B (GPTQ-4) & GSM8K & 1/10 & 138 [138, 140] & 2996 [2897, 3103] & 41.16 [39.95, 42.22] & 116.4 \\
GPU server & Phi-4 Mini 3.8B (GPTQ-4) & HellaSwag & 1/10 & 139 [138, 140] & 253 [251, 255] & 30.61 [30.05, 31.17] & 3.5 \\
GPU server & Phi-4 Mini 3.8B (GPTQ-4) & Dolly & 1/10 & 140 [139, 141] & 2498 [2380, 2626] & 42.50 [41.41, 43.52] & 103.9 \\
GPU server & Phi-4 Mini 3.8B (bf16) & MMLU & 1/10 & 138 [137, 139] & 251 [250, 253] & 23.89 [23.26, 24.49] & 2.7 \\
GPU server & Phi-4 Mini 3.8B (bf16) & GSM8K & 1/10 & 136 [135, 137] & 2918 [2837, 2997] & 42.99 [42.12, 43.85] & 117.7 \\
GPU server & Phi-4 Mini 3.8B (bf16) & HellaSwag & 1/10 & 137 [136, 138] & 249 [247, 251] & 31.67 [31.11, 32.21] & 3.6 \\
GPU server & Phi-4 Mini 3.8B (bf16) & Dolly & 1/10 & 140 [139, 142] & 2545 [2432, 2652] & 44.18 [42.94, 45.46] & 110.4 \\
GPU server & Qwen 2.5 1.5B (bf16) & MMLU & 1/10 & 132 [130, 133] & 242 [240, 245] & 9.23 [9.10, 9.38] & 1.0 \\
GPU server & Qwen 2.5 1.5B (bf16) & GSM8K & 1/10 & 130 [129, 130] & 2455 [2407, 2497] & 62.02 [60.93, 63.30] & 144.0 \\
GPU server & Qwen 2.5 1.5B (bf16) & HellaSwag & 1/10 & 134 [133, 135] & 246 [244, 248] & 8.94 [8.87, 9.02] & 1.0 \\
GPU server & Qwen 2.5 1.5B (bf16) & Dolly & 1/10 & 132 [131, 133] & 1604 [1507, 1702] & 69.21 [66.60, 71.75] & 107.7 \\
GPU server & Qwen 2.5 7B (AWQ) & MMLU & 1/10 & 135 [134, 136] & 247 [244, 249] & 9.00 [8.90, 9.09] & 1.0 \\
GPU server & Qwen 2.5 7B (AWQ) & GSM8K & 1/10 & 132 [130, 133] & 2465 [2411, 2530] & 59.22 [57.79, 60.51] & 137.3 \\
GPU server & Qwen 2.5 7B (AWQ) & HellaSwag & 1/10 & 137 [136, 137] & 248 [247, 250] & 8.97 [8.90, 9.04] & 1.0 \\
GPU server & Qwen 2.5 7B (AWQ) & Dolly & 1/10 & 135 [134, 136] & 1974 [1874, 2074] & 67.84 [66.34, 69.39] & 128.0 \\
GPU server & Qwen 3 30B-A3B (bf16) & MMLU & 1/10 & 152 [150, 153] & 265 [263, 268] & 8.83 [8.74, 8.90] & 1.0 \\
GPU server & Qwen 3 30B-A3B (bf16) & GSM8K & 1/10 & 152 [151, 154] & 6935 [6749, 7128] & 16.60 [16.22, 16.99] & 112.8 \\
GPU server & Qwen 3 30B-A3B (bf16) & HellaSwag & 1/10 & 151 [150, 152] & 264 [263, 266] & 8.85 [8.80, 8.90] & 1.0 \\
GPU server & Qwen 3 30B-A3B (bf16) & Dolly & 1/10 & 155 [154, 156] & 6103 [5840, 6354] & 19.72 [19.20, 20.22] & 120.1 \\
GPU server & Qwen 3 4B (AWQ) & MMLU & 1/10 & 135 [134, 136] & 247 [246, 249] & 10.18 [9.89, 10.46] & 1.1 \\
GPU server & Qwen 3 4B (AWQ) & GSM8K & 1/10 & 134 [134, 135] & 3153 [3070, 3242] & 43.38 [42.21, 44.50] & 131.1 \\
GPU server & Qwen 3 4B (AWQ) & HellaSwag & 1/10 & 137 [136, 139] & 250 [248, 252] & 8.87 [8.79, 8.95] & 1.0 \\
GPU server & Qwen 3 4B (AWQ) & Dolly & 1/10 & 135 [135, 136] & 2807 [2698, 2910] & 48.46 [47.22, 49.74] & 132.4 \\
\end{longtable}

\begin{longtable}{@{}l l l r r@{}}
\caption{Run-level request rate and sustained throughput (mean over runs with 95\% $t$-intervals; sessions weighted equally). Both include all inter-request overheads and, on the Raspberry Pi~5, the cooling pause (1\,s for MMLU and HellaSwag, 2\,s for GSM8K and Dolly). One sequential request stream.}\label{tab:s2}\\
\toprule
Context & Configuration & Workload & Requests/s & Sustained words/s \\
\midrule\endfirsthead
\toprule
Context & Configuration & Workload & Requests/s & Sustained words/s \\
\midrule\endhead
\bottomrule\endfoot
Edge & Gemma 3 1B (Q8) & MMLU & 0.391 [0.389, 0.392] & 0.39 [0.39, 0.39] \\
Edge & Gemma 3 1B (Q8) & GSM8K & 0.039 [0.039, 0.039] & 5.47 [5.45, 5.48] \\
Edge & Gemma 3 1B (Q8) & HellaSwag & 0.328 [0.326, 0.330] & 0.33 [0.33, 0.33] \\
Edge & Llama 3.2 1B (Q4) & MMLU & 0.361 [0.361, 0.362] & 0.36 [0.36, 0.36] \\
Edge & Llama 3.2 1B (Q4) & GSM8K & 0.056 [0.055, 0.056] & 8.01 [8.00, 8.03] \\
Edge & Llama 3.2 1B (Q4) & HellaSwag & 0.292 [0.292, 0.292] & 0.29 [0.29, 0.29] \\
Edge & Phi-4 Mini 3.8B (Q4) & MMLU & 0.133 [0.132, 0.133] & 0.34 [0.34, 0.34] \\
Edge & Phi-4 Mini 3.8B (Q4) & GSM8K & 0.023 [0.023, 0.023] & 2.73 [2.72, 2.75] \\
Edge & Phi-4 Mini 3.8B (Q4) & HellaSwag & 0.102 [0.101, 0.102] & 0.30 [0.30, 0.30] \\
Edge & SmolLM3 3B (Q4) & MMLU & 0.150 [0.149, 0.151] & 0.42 [0.41, 0.42] \\
Edge & SmolLM3 3B (Q4) & GSM8K & 0.024 [0.024, 0.024] & 3.58 [3.56, 3.60] \\
Edge & SmolLM3 3B (Q4) & HellaSwag & 0.118 [0.118, 0.119] & 0.36 [0.36, 0.37] \\
\midrule
Local CPU & Gemma 3 1B (Q8) & MMLU & 0.863 [0.844, 0.881] & 0.86 [0.84, 0.88] \\
Local CPU & Gemma 3 1B (Q8) & GSM8K & 0.095 [0.094, 0.095] & 13.22 [13.17, 13.26] \\
Local CPU & Gemma 3 1B (Q8) & HellaSwag & 0.633 [0.625, 0.640] & 0.63 [0.63, 0.64] \\
Local CPU & Llama 3.2 1B (Q8) & MMLU & 0.538 [0.528, 0.548] & 0.54 [0.53, 0.55] \\
Local CPU & Llama 3.2 1B (Q8) & GSM8K & 0.307 [0.305, 0.309] & 8.58 [8.52, 8.64] \\
Local CPU & Llama 3.2 1B (Q8) & HellaSwag & 0.400 [0.397, 0.402] & 0.40 [0.40, 0.40] \\
Local CPU & Ministral 3 3B (Q4) & MMLU & 0.177 [0.176, 0.177] & 0.18 [0.18, 0.18] \\
Local CPU & Ministral 3 3B (Q4) & GSM8K & 0.043 [0.043, 0.044] & 5.06 [5.04, 5.09] \\
Local CPU & Ministral 3 3B (Q4) & HellaSwag & 0.131 [0.131, 0.132] & 0.13 [0.13, 0.13] \\
Local CPU & Phi-4 Mini 3.8B (Q4) & MMLU & 0.229 [0.227, 0.232] & 0.58 [0.57, 0.58] \\
Local CPU & Phi-4 Mini 3.8B (Q4) & GSM8K & 0.054 [0.053, 0.055] & 6.32 [6.21, 6.43] \\
Local CPU & Phi-4 Mini 3.8B (Q4) & HellaSwag & 0.157 [0.157, 0.158] & 0.47 [0.47, 0.47] \\
\midrule
GPU server & Apertus 70B (bf16) & MMLU & 2.203 [2.197, 2.210] & 2.27 [2.26, 2.28] \\
GPU server & Apertus 70B (bf16) & GSM8K & 0.090 [0.090, 0.090] & 11.67 [11.65, 11.70] \\
GPU server & Apertus 70B (bf16) & HellaSwag & 2.179 [2.173, 2.185] & 2.18 [2.17, 2.19] \\
GPU server & Apertus 70B (bf16) & Dolly & 0.094 [0.092, 0.097] & 12.86 [12.72, 12.99] \\
GPU server & Granite 4 H-Tiny (bf16) & MMLU & 2.463 [2.455, 2.471] & 2.68 [2.68, 2.69] \\
GPU server & Granite 4 H-Tiny (bf16) & GSM8K & 0.135 [0.132, 0.137] & 14.29 [14.02, 14.56] \\
GPU server & Granite 4 H-Tiny (bf16) & HellaSwag & 2.456 [2.443, 2.469] & 2.73 [2.71, 2.74] \\
GPU server & Granite 4 H-Tiny (bf16) & Dolly & 0.135 [0.130, 0.140] & 15.73 [15.42, 16.05] \\
GPU server & Kimi Linear 48B-A3B (bf16) & MMLU & 2.350 [2.329, 2.371] & 4.82 [4.77, 4.86] \\
GPU server & Kimi Linear 48B-A3B (bf16) & GSM8K & 0.167 [0.161, 0.174] & 14.97 [14.37, 15.57] \\
GPU server & Kimi Linear 48B-A3B (bf16) & HellaSwag & 2.388 [2.374, 2.402] & 4.32 [4.30, 4.35] \\
GPU server & Kimi Linear 48B-A3B (bf16) & Dolly & 0.183 [0.174, 0.192] & 16.73 [16.05, 17.40] \\
GPU server & Llama 3.2 1B (bf16) & MMLU & 2.862 [2.827, 2.898] & 2.86 [2.83, 2.90] \\
GPU server & Llama 3.2 1B (bf16) & GSM8K & 0.828 [0.812, 0.845] & 93.24 [91.41, 95.06] \\
GPU server & Llama 3.2 1B (bf16) & HellaSwag & 2.844 [2.814, 2.875] & 2.84 [2.81, 2.88] \\
GPU server & Llama 3.2 1B (bf16) & Dolly & 0.814 [0.789, 0.839] & 98.17 [96.02, 100.32] \\
GPU server & Llama 3.3 70B (FP8) & MMLU & 2.469 [2.457, 2.482] & 2.47 [2.46, 2.48] \\
GPU server & Llama 3.3 70B (FP8) & GSM8K & 0.141 [0.141, 0.141] & 20.14 [20.13, 20.14] \\
GPU server & Llama 3.3 70B (FP8) & HellaSwag & 2.406 [2.393, 2.420] & 2.41 [2.39, 2.42] \\
GPU server & Llama 3.3 70B (FP8) & Dolly & 0.157 [0.151, 0.162] & 18.06 [17.81, 18.31] \\
GPU server & Ministral 3 14B & MMLU & 2.597 [2.572, 2.623] & 2.60 [2.57, 2.62] \\
GPU server & Ministral 3 14B & GSM8K & 0.086 [0.085, 0.087] & 10.49 [10.34, 10.64] \\
GPU server & Ministral 3 14B & HellaSwag & 2.589 [2.572, 2.607] & 2.59 [2.57, 2.61] \\
GPU server & Ministral 3 14B & Dolly & 0.100 [0.094, 0.106] & 11.54 [10.96, 12.13] \\
GPU server & Ministral 3 3B (AWQ) & MMLU & 2.839 [2.805, 2.873] & 2.84 [2.80, 2.87] \\
GPU server & Ministral 3 3B (AWQ) & GSM8K & 0.331 [0.319, 0.343] & 41.40 [39.88, 42.91] \\
GPU server & Ministral 3 3B (AWQ) & HellaSwag & 2.825 [2.795, 2.855] & 2.82 [2.79, 2.85] \\
GPU server & Ministral 3 3B (AWQ) & Dolly & 0.379 [0.368, 0.390] & 45.63 [44.21, 47.05] \\
GPU server & Ministral 3 8B (AWQ) & MMLU & 2.827 [2.807, 2.846] & 2.83 [2.81, 2.85] \\
GPU server & Ministral 3 8B (AWQ) & GSM8K & 0.274 [0.269, 0.280] & 34.73 [33.99, 35.47] \\
GPU server & Ministral 3 8B (AWQ) & HellaSwag & 2.776 [2.752, 2.799] & 2.78 [2.75, 2.80] \\
GPU server & Ministral 3 8B (AWQ) & Dolly & 0.302 [0.299, 0.305] & 36.33 [35.00, 37.65] \\
GPU server & Phi-4 14B (bf16) & MMLU & 2.690 [2.661, 2.719] & 3.50 [3.46, 3.53] \\
GPU server & Phi-4 14B (bf16) & GSM8K & 0.337 [0.333, 0.340] & 45.46 [44.97, 45.95] \\
GPU server & Phi-4 14B (bf16) & HellaSwag & 2.671 [2.649, 2.693] & 2.99 [2.97, 3.02] \\
GPU server & Phi-4 14B (bf16) & Dolly & 0.387 [0.375, 0.399] & 51.25 [50.79, 51.71] \\
GPU server & Phi-4 Mini 3.8B (GPTQ-4) & MMLU & 2.748 [2.735, 2.762] & 6.40 [6.37, 6.44] \\
GPU server & Phi-4 Mini 3.8B (GPTQ-4) & GSM8K & 0.322 [0.312, 0.332] & 37.50 [36.37, 38.64] \\
GPU server & Phi-4 Mini 3.8B (GPTQ-4) & HellaSwag & 2.747 [2.719, 2.776] & 9.53 [9.43, 9.63] \\
GPU server & Phi-4 Mini 3.8B (GPTQ-4) & Dolly & 0.384 [0.371, 0.397] & 39.82 [38.92, 40.71] \\
GPU server & Phi-4 Mini 3.8B (bf16) & MMLU & 2.770 [2.747, 2.793] & 7.48 [7.42, 7.54] \\
GPU server & Phi-4 Mini 3.8B (bf16) & GSM8K & 0.330 [0.323, 0.338] & 38.89 [38.04, 39.74] \\
GPU server & Phi-4 Mini 3.8B (bf16) & HellaSwag & 2.789 [2.761, 2.817] & 9.93 [9.83, 10.03] \\
GPU server & Phi-4 Mini 3.8B (bf16) & Dolly & 0.377 [0.365, 0.389] & 41.58 [40.29, 42.86] \\
GPU server & Qwen 2.5 1.5B (bf16) & MMLU & 2.838 [2.801, 2.874] & 2.89 [2.86, 2.93] \\
GPU server & Qwen 2.5 1.5B (bf16) & GSM8K & 0.390 [0.383, 0.397] & 56.18 [55.12, 57.24] \\
GPU server & Qwen 2.5 1.5B (bf16) & HellaSwag & 2.795 [2.767, 2.824] & 2.80 [2.77, 2.82] \\
GPU server & Qwen 2.5 1.5B (bf16) & Dolly & 0.586 [0.553, 0.619] & 62.88 [60.37, 65.39] \\
GPU server & Qwen 2.5 7B (AWQ) & MMLU & 2.803 [2.771, 2.836] & 2.80 [2.77, 2.84] \\
GPU server & Qwen 2.5 7B (AWQ) & GSM8K & 0.389 [0.379, 0.398] & 53.37 [52.05, 54.70] \\
GPU server & Qwen 2.5 7B (AWQ) & HellaSwag & 2.792 [2.771, 2.814] & 2.79 [2.77, 2.81] \\
GPU server & Qwen 2.5 7B (AWQ) & Dolly & 0.481 [0.458, 0.505] & 61.44 [59.93, 62.95] \\
GPU server & Qwen 3 30B-A3B (bf16) & MMLU & 2.660 [2.632, 2.689] & 2.66 [2.63, 2.69] \\
GPU server & Qwen 3 30B-A3B (bf16) & GSM8K & 0.142 [0.138, 0.146] & 16.03 [15.60, 16.45] \\
GPU server & Qwen 3 30B-A3B (bf16) & HellaSwag & 2.670 [2.654, 2.685] & 2.67 [2.65, 2.69] \\
GPU server & Qwen 3 30B-A3B (bf16) & Dolly & 0.161 [0.156, 0.167] & 19.32 [18.84, 19.80] \\
GPU server & Qwen 3 4B (AWQ) & MMLU & 2.797 [2.777, 2.817] & 3.19 [3.17, 3.21] \\
GPU server & Qwen 3 4B (AWQ) & GSM8K & 0.307 [0.298, 0.315] & 40.21 [39.07, 41.34] \\
GPU server & Qwen 3 4B (AWQ) & HellaSwag & 2.766 [2.739, 2.793] & 2.77 [2.74, 2.79] \\
GPU server & Qwen 3 4B (AWQ) & Dolly & 0.343 [0.332, 0.355] & 45.42 [44.09, 46.75] \\
\end{longtable}

\endgroup

\begingroup\small
\begin{longtable}{@{}l l r r@{}}
\caption{Mean GPU energy per request on the GPU server with 95\% bootstrap intervals over requests. Energy is the trapezoidal integral of the board power of the first GPU reported by \texttt{nvidia-smi} over the request; all models were served with tensor parallelism across both GPUs, so values cover one of the two GPUs.}\label{tab:s3}\\
\toprule
Configuration & Workload & Requests & Energy (J) \\
\midrule\endfirsthead
\toprule
Configuration & Workload & Requests & Energy (J) \\
\midrule\endhead
\bottomrule\endfoot
Apertus 70B (bf16) & MMLU & 100 & 103 [97, 108] \\
Apertus 70B (bf16) & GSM8K & 100 & 5527 [5499, 5550] \\
Apertus 70B (bf16) & HellaSwag & 100 & 92 [90, 93] \\
Apertus 70B (bf16) & Dolly & 100 & 5603 [5591, 5617] \\
Granite 4 H-Tiny (bf16) & MMLU & 200 & 33 [33, 34] \\
Granite 4 H-Tiny (bf16) & GSM8K & 200 & 1252 [1228, 1274] \\
Granite 4 H-Tiny (bf16) & HellaSwag & 200 & 34 [34, 35] \\
Granite 4 H-Tiny (bf16) & Dolly & 100 & 1333 [1325, 1341] \\
Kimi Linear 48B-A3B (bf16) & MMLU & 100 & 40 [40, 41] \\
Kimi Linear 48B-A3B (bf16) & GSM8K & 100 & 1202 [1181, 1221] \\
Kimi Linear 48B-A3B (bf16) & HellaSwag & 100 & 41 [41, 41] \\
Kimi Linear 48B-A3B (bf16) & Dolly & 100 & 1184 [1168, 1197] \\
Llama 3.2 1B (bf16) & MMLU & 100 & 23 [23, 24] \\
Llama 3.2 1B (bf16) & GSM8K & 100 & 208 [203, 213] \\
Llama 3.2 1B (bf16) & HellaSwag & 100 & 24 [23, 24] \\
Llama 3.2 1B (bf16) & Dolly & 100 & 226 [225, 227] \\
Llama 3.3 70B (FP8) & MMLU & 200 & 59 [58, 61] \\
Llama 3.3 70B (FP8) & GSM8K & 200 & 2875 [2832, 2920] \\
Llama 3.3 70B (FP8) & HellaSwag & 200 & 58 [57, 58] \\
Llama 3.3 70B (FP8) & Dolly & 100 & 3183 [3153, 3199] \\
Ministral 3 14B & MMLU & 100 & 35 [35, 36] \\
Ministral 3 14B & GSM8K & 100 & 1981 [1974, 1989] \\
Ministral 3 14B & HellaSwag & 100 & 35 [35, 36] \\
Ministral 3 14B & Dolly & 100 & 2151 [2130, 2171] \\
Ministral 3 3B (AWQ) & MMLU & 100 & 24 [24, 25] \\
Ministral 3 3B (AWQ) & GSM8K & 100 & 395 [394, 396] \\
Ministral 3 3B (AWQ) & HellaSwag & 100 & 25 [24, 25] \\
Ministral 3 3B (AWQ) & Dolly & 100 & 390 [389, 392] \\
Ministral 3 8B (AWQ) & MMLU & 100 & 27 [26, 27] \\
Ministral 3 8B (AWQ) & GSM8K & 100 & 540 [537, 544] \\
Ministral 3 8B (AWQ) & HellaSwag & 100 & 27 [27, 27] \\
Ministral 3 8B (AWQ) & Dolly & 100 & 546 [544, 547] \\
Phi-4 14B (bf16) & MMLU & 100 & 37 [37, 38] \\
Phi-4 14B (bf16) & GSM8K & 100 & 1214 [1206, 1222] \\
Phi-4 14B (bf16) & HellaSwag & 100 & 37 [36, 37] \\
Phi-4 14B (bf16) & Dolly & 100 & 1209 [1199, 1216] \\
Phi-4 Mini 3.8B (GPTQ-4) & MMLU & 100 & 26 [26, 26] \\
Phi-4 Mini 3.8B (GPTQ-4) & GSM8K & 100 & 523 [515, 531] \\
Phi-4 Mini 3.8B (GPTQ-4) & HellaSwag & 100 & 27 [27, 27] \\
Phi-4 Mini 3.8B (GPTQ-4) & Dolly & 100 & 532 [522, 540] \\
Phi-4 Mini 3.8B (bf16) & MMLU & 100 & 29 [28, 29] \\
Phi-4 Mini 3.8B (bf16) & GSM8K & 100 & 635 [626, 643] \\
Phi-4 Mini 3.8B (bf16) & HellaSwag & 100 & 28 [28, 28] \\
Phi-4 Mini 3.8B (bf16) & Dolly & 100 & 659 [649, 668] \\
Qwen 2.5 1.5B (bf16) & MMLU & 100 & 24 [24, 24] \\
Qwen 2.5 1.5B (bf16) & GSM8K & 100 & 357 [356, 358] \\
Qwen 2.5 1.5B (bf16) & HellaSwag & 100 & 24 [24, 24] \\
Qwen 2.5 1.5B (bf16) & Dolly & 100 & 359 [353, 364] \\
Qwen 2.5 7B (AWQ) & MMLU & 100 & 26 [26, 27] \\
Qwen 2.5 7B (AWQ) & GSM8K & 100 & 407 [404, 410] \\
Qwen 2.5 7B (AWQ) & HellaSwag & 100 & 26 [26, 27] \\
Qwen 2.5 7B (AWQ) & Dolly & 100 & 405 [403, 407] \\
Qwen 3 30B-A3B (bf16) & MMLU & 100 & 31 [30, 31] \\
Qwen 3 30B-A3B (bf16) & GSM8K & 100 & 1033 [1026, 1039] \\
Qwen 3 30B-A3B (bf16) & HellaSwag & 100 & 31 [31, 31] \\
Qwen 3 30B-A3B (bf16) & Dolly & 100 & 1041 [1028, 1052] \\
Qwen 3 4B (AWQ) & MMLU & 100 & 25 [25, 26] \\
Qwen 3 4B (AWQ) & GSM8K & 100 & 451 [448, 453] \\
Qwen 3 4B (AWQ) & HellaSwag & 100 & 26 [25, 26] \\
Qwen 3 4B (AWQ) & Dolly & 100 & 447 [445, 449] \\
\end{longtable}

\endgroup

\begingroup\scriptsize\setlength{\tabcolsep}{4pt}
\begin{longtable}{@{}l l l r r r r@{}}
\caption{Run-to-run spread of accuracy. Accuracy is the session-weighted mean; the spread is computed over all measured runs. Because the same 100 items are used in every run with greedy decoding, the spread reflects execution non-determinism, not item sampling; item-level uncertainty is given by the Wilson intervals in the article.}\label{tab:s4}\\
\toprule
Context & Configuration & Workload & Accuracy & Min & Max & SD \\
\midrule\endfirsthead
\toprule
Context & Configuration & Workload & Accuracy & Min & Max & SD \\
\midrule\endhead
\bottomrule\endfoot
Edge & Gemma 3 1B (Q8) & MMLU & 0.315 & 0.30 & 0.34 & 0.013 \\
Edge & Gemma 3 1B (Q8) & GSM8K & 0.477 & 0.45 & 0.51 & 0.022 \\
Edge & Gemma 3 1B (Q8) & HellaSwag & 0.291 & 0.28 & 0.31 & 0.012 \\
Edge & Llama 3.2 1B (Q4) & MMLU & 0.395 & 0.38 & 0.41 & 0.011 \\
Edge & Llama 3.2 1B (Q4) & GSM8K & 0.520 & 0.49 & 0.55 & 0.017 \\
Edge & Llama 3.2 1B (Q4) & HellaSwag & 0.251 & 0.23 & 0.26 & 0.010 \\
Edge & Phi-4 Mini 3.8B (Q4) & MMLU & 0.564 & 0.54 & 0.58 & 0.012 \\
Edge & Phi-4 Mini 3.8B (Q4) & GSM8K & 0.639 & 0.60 & 0.66 & 0.018 \\
Edge & Phi-4 Mini 3.8B (Q4) & HellaSwag & 0.489 & 0.48 & 0.50 & 0.009 \\
Edge & SmolLM3 3B (Q4) & MMLU & 0.520 & 0.51 & 0.53 & 0.005 \\
Edge & SmolLM3 3B (Q4) & GSM8K & 0.577 & 0.56 & 0.62 & 0.018 \\
Edge & SmolLM3 3B (Q4) & HellaSwag & 0.290 & 0.28 & 0.30 & 0.008 \\
\midrule
Local CPU & Gemma 3 1B (Q8) & MMLU & 0.310 & 0.31 & 0.31 & 0.000 \\
Local CPU & Gemma 3 1B (Q8) & GSM8K & 0.471 & 0.44 & 0.50 & 0.019 \\
Local CPU & Gemma 3 1B (Q8) & HellaSwag & 0.290 & 0.29 & 0.29 & 0.000 \\
Local CPU & Llama 3.2 1B (Q8) & MMLU & 0.380 & 0.38 & 0.38 & 0.000 \\
Local CPU & Llama 3.2 1B (Q8) & GSM8K & 0.140 & 0.14 & 0.14 & 0.000 \\
Local CPU & Llama 3.2 1B (Q8) & HellaSwag & 0.290 & 0.29 & 0.29 & 0.000 \\
Local CPU & Ministral 3 3B (Q4) & MMLU & 0.630 & 0.62 & 0.64 & 0.010 \\
Local CPU & Ministral 3 3B (Q4) & GSM8K & 0.549 & 0.51 & 0.58 & 0.019 \\
Local CPU & Ministral 3 3B (Q4) & HellaSwag & 0.405 & 0.39 & 0.42 & 0.015 \\
Local CPU & Phi-4 Mini 3.8B (Q4) & MMLU & 0.540 & 0.54 & 0.54 & 0.000 \\
Local CPU & Phi-4 Mini 3.8B (Q4) & GSM8K & 0.640 & 0.64 & 0.64 & 0.000 \\
Local CPU & Phi-4 Mini 3.8B (Q4) & HellaSwag & 0.510 & 0.50 & 0.52 & 0.005 \\
\midrule
GPU server & Apertus 70B (bf16) & MMLU & 0.610 & 0.61 & 0.61 & 0.000 \\
GPU server & Apertus 70B (bf16) & GSM8K & 0.601 & 0.55 & 0.61 & 0.020 \\
GPU server & Apertus 70B (bf16) & HellaSwag & 0.410 & 0.41 & 0.41 & 0.000 \\
GPU server & Granite 4 H-Tiny (bf16) & MMLU & 0.580 & 0.58 & 0.58 & 0.000 \\
GPU server & Granite 4 H-Tiny (bf16) & GSM8K & 0.690 & 0.69 & 0.69 & 0.000 \\
GPU server & Granite 4 H-Tiny (bf16) & HellaSwag & 0.640 & 0.64 & 0.64 & 0.000 \\
GPU server & Kimi Linear 48B-A3B (bf16) & MMLU & 0.450 & 0.45 & 0.45 & 0.000 \\
GPU server & Kimi Linear 48B-A3B (bf16) & GSM8K & 0.710 & 0.71 & 0.71 & 0.000 \\
GPU server & Kimi Linear 48B-A3B (bf16) & HellaSwag & 0.240 & 0.24 & 0.24 & 0.000 \\
GPU server & Llama 3.2 1B (bf16) & MMLU & 0.430 & 0.43 & 0.43 & 0.000 \\
GPU server & Llama 3.2 1B (bf16) & GSM8K & 0.435 & 0.43 & 0.44 & 0.005 \\
GPU server & Llama 3.2 1B (bf16) & HellaSwag & 0.230 & 0.23 & 0.23 & 0.000 \\
GPU server & Llama 3.3 70B (FP8) & MMLU & 0.802 & 0.80 & 0.81 & 0.004 \\
GPU server & Llama 3.3 70B (FP8) & GSM8K & 0.745 & 0.74 & 0.75 & 0.005 \\
GPU server & Llama 3.3 70B (FP8) & HellaSwag & 0.835 & 0.83 & 0.84 & 0.005 \\
GPU server & Ministral 3 14B & MMLU & 0.700 & 0.70 & 0.70 & 0.000 \\
GPU server & Ministral 3 14B & GSM8K & 0.590 & 0.59 & 0.59 & 0.000 \\
GPU server & Ministral 3 14B & HellaSwag & 0.570 & 0.57 & 0.57 & 0.000 \\
GPU server & Ministral 3 3B (AWQ) & MMLU & 0.640 & 0.64 & 0.64 & 0.000 \\
GPU server & Ministral 3 3B (AWQ) & GSM8K & 0.510 & 0.51 & 0.51 & 0.000 \\
GPU server & Ministral 3 3B (AWQ) & HellaSwag & 0.460 & 0.46 & 0.46 & 0.000 \\
GPU server & Ministral 3 8B (AWQ) & MMLU & 0.660 & 0.66 & 0.66 & 0.000 \\
GPU server & Ministral 3 8B (AWQ) & GSM8K & 0.510 & 0.51 & 0.51 & 0.000 \\
GPU server & Ministral 3 8B (AWQ) & HellaSwag & 0.640 & 0.64 & 0.64 & 0.000 \\
GPU server & Phi-4 14B (bf16) & MMLU & 0.820 & 0.82 & 0.82 & 0.000 \\
GPU server & Phi-4 14B (bf16) & GSM8K & 0.550 & 0.55 & 0.55 & 0.000 \\
GPU server & Phi-4 14B (bf16) & HellaSwag & 0.850 & 0.85 & 0.85 & 0.000 \\
GPU server & Phi-4 Mini 3.8B (GPTQ-4) & MMLU & 0.600 & 0.60 & 0.60 & 0.000 \\
GPU server & Phi-4 Mini 3.8B (GPTQ-4) & GSM8K & 0.610 & 0.61 & 0.61 & 0.000 \\
GPU server & Phi-4 Mini 3.8B (GPTQ-4) & HellaSwag & 0.550 & 0.55 & 0.55 & 0.000 \\
GPU server & Phi-4 Mini 3.8B (bf16) & MMLU & 0.630 & 0.63 & 0.63 & 0.000 \\
GPU server & Phi-4 Mini 3.8B (bf16) & GSM8K & 0.670 & 0.67 & 0.67 & 0.000 \\
GPU server & Phi-4 Mini 3.8B (bf16) & HellaSwag & 0.530 & 0.53 & 0.53 & 0.000 \\
GPU server & Qwen 2.5 1.5B (bf16) & MMLU & 0.520 & 0.52 & 0.52 & 0.000 \\
GPU server & Qwen 2.5 1.5B (bf16) & GSM8K & 0.390 & 0.39 & 0.39 & 0.000 \\
GPU server & Qwen 2.5 1.5B (bf16) & HellaSwag & 0.250 & 0.25 & 0.25 & 0.000 \\
GPU server & Qwen 2.5 7B (AWQ) & MMLU & 0.670 & 0.67 & 0.67 & 0.000 \\
GPU server & Qwen 2.5 7B (AWQ) & GSM8K & 0.550 & 0.55 & 0.55 & 0.000 \\
GPU server & Qwen 2.5 7B (AWQ) & HellaSwag & 0.590 & 0.59 & 0.59 & 0.000 \\
GPU server & Qwen 3 30B-A3B (bf16) & MMLU & 0.816 & 0.81 & 0.82 & 0.005 \\
GPU server & Qwen 3 30B-A3B (bf16) & GSM8K & 0.590 & 0.59 & 0.59 & 0.000 \\
GPU server & Qwen 3 30B-A3B (bf16) & HellaSwag & 0.790 & 0.79 & 0.79 & 0.000 \\
GPU server & Qwen 3 4B (AWQ) & MMLU & 0.710 & 0.71 & 0.71 & 0.000 \\
GPU server & Qwen 3 4B (AWQ) & GSM8K & 0.540 & 0.54 & 0.54 & 0.000 \\
GPU server & Qwen 3 4B (AWQ) & HellaSwag & 0.690 & 0.69 & 0.69 & 0.000 \\
\end{longtable}

\begin{longtable}{@{}l l l l r r r c@{}}
\caption{Paired sensitivity analysis: exact McNemar tests of the configuration with the highest observed accuracy against every other configuration in the same context and benchmark, on the identical 100 items (item correct if correct in at least half of its runs), Holm-adjusted within each cell. L/O: items answered correctly only by the leader / only by the other configuration.}\label{tab:s5}\\
\toprule
Context & Bench. & Leader & Compared configuration & L/O & $p$ & $p_{\mathrm{Holm}}$ & Resolved \\
\midrule\endfirsthead
\toprule
Context & Bench. & Leader & Compared configuration & L/O & $p$ & $p_{\mathrm{Holm}}$ & Resolved \\
\midrule\endhead
\bottomrule\endfoot
Edge & MMLU & Phi-4 Mini 3.8B (Q4) & Gemma 3 1B (Q8) & 34/8 & $<$0.001 & $<$0.001 & yes \\
Edge & MMLU & Phi-4 Mini 3.8B (Q4) & Llama 3.2 1B (Q4) & 28/11 & 0.009 & 0.019 & yes \\
Edge & MMLU & Phi-4 Mini 3.8B (Q4) & SmolLM3 3B (Q4) & 15/10 & 0.424 & 0.424 & no \\
\midrule
Edge & GSM8K & Phi-4 Mini 3.8B (Q4) & Gemma 3 1B (Q8) & 26/8 & 0.003 & 0.009 & yes \\
Edge & GSM8K & Phi-4 Mini 3.8B (Q4) & Llama 3.2 1B (Q4) & 22/9 & 0.029 & 0.059 & no \\
Edge & GSM8K & Phi-4 Mini 3.8B (Q4) & SmolLM3 3B (Q4) & 11/5 & 0.210 & 0.210 & no \\
\midrule
Edge & HellaSwag & Phi-4 Mini 3.8B (Q4) & Llama 3.2 1B (Q4) & 34/10 & $<$0.001 & 0.001 & yes \\
Edge & HellaSwag & Phi-4 Mini 3.8B (Q4) & SmolLM3 3B (Q4) & 32/12 & 0.004 & 0.007 & yes \\
Edge & HellaSwag & Phi-4 Mini 3.8B (Q4) & Gemma 3 1B (Q8) & 38/20 & 0.025 & 0.025 & yes \\
\midrule
Local CPU & MMLU & Ministral 3 3B (Q4) & Gemma 3 1B (Q8) & 41/4 & $<$0.001 & $<$0.001 & yes \\
Local CPU & MMLU & Ministral 3 3B (Q4) & Llama 3.2 1B (Q8) & 38/8 & $<$0.001 & $<$0.001 & yes \\
Local CPU & MMLU & Ministral 3 3B (Q4) & Phi-4 Mini 3.8B (Q4) & 20/6 & 0.009 & 0.009 & yes \\
\midrule
Local CPU & GSM8K & Phi-4 Mini 3.8B (Q4) & Llama 3.2 1B (Q8) & 51/1 & $<$0.001 & $<$0.001 & yes \\
Local CPU & GSM8K & Phi-4 Mini 3.8B (Q4) & Gemma 3 1B (Q8) & 27/10 & 0.008 & 0.015 & yes \\
Local CPU & GSM8K & Phi-4 Mini 3.8B (Q4) & Ministral 3 3B (Q4) & 18/6 & 0.023 & 0.023 & yes \\
\midrule
Local CPU & HellaSwag & Phi-4 Mini 3.8B (Q4) & Llama 3.2 1B (Q8) & 31/9 & $<$0.001 & 0.002 & yes \\
Local CPU & HellaSwag & Phi-4 Mini 3.8B (Q4) & Gemma 3 1B (Q8) & 40/18 & 0.005 & 0.011 & yes \\
Local CPU & HellaSwag & Phi-4 Mini 3.8B (Q4) & Ministral 3 3B (Q4) & 26/20 & 0.461 & 0.461 & no \\
\midrule
GPU server & MMLU & Phi-4 14B (bf16) & Llama 3.2 1B (bf16) & 42/3 & $<$0.001 & $<$0.001 & yes \\
GPU server & MMLU & Phi-4 14B (bf16) & Kimi Linear 48B-A3B (bf16) & 40/3 & $<$0.001 & $<$0.001 & yes \\
GPU server & MMLU & Phi-4 14B (bf16) & Qwen 2.5 1.5B (bf16) & 36/6 & $<$0.001 & $<$0.001 & yes \\
GPU server & MMLU & Phi-4 14B (bf16) & Granite 4 H-Tiny (bf16) & 27/3 & $<$0.001 & $<$0.001 & yes \\
GPU server & MMLU & Phi-4 14B (bf16) & Apertus 70B (bf16) & 26/5 & $<$0.001 & 0.002 & yes \\
GPU server & MMLU & Phi-4 14B (bf16) & Phi-4 Mini 3.8B (GPTQ-4) & 28/6 & $<$0.001 & 0.002 & yes \\
GPU server & MMLU & Phi-4 14B (bf16) & Phi-4 Mini 3.8B (bf16) & 23/4 & $<$0.001 & 0.002 & yes \\
GPU server & MMLU & Phi-4 14B (bf16) & Qwen 2.5 7B (AWQ) & 18/3 & 0.001 & 0.010 & yes \\
GPU server & MMLU & Phi-4 14B (bf16) & Ministral 3 3B (AWQ) & 25/7 & 0.002 & 0.013 & yes \\
GPU server & MMLU & Phi-4 14B (bf16) & Ministral 3 8B (AWQ) & 21/5 & 0.002 & 0.013 & yes \\
GPU server & MMLU & Phi-4 14B (bf16) & Ministral 3 14B & 20/8 & 0.036 & 0.143 & no \\
GPU server & MMLU & Phi-4 14B (bf16) & Qwen 3 4B (AWQ) & 18/7 & 0.043 & 0.143 & no \\
GPU server & MMLU & Phi-4 14B (bf16) & Llama 3.3 70B (FP8) & 11/10 & 1.000 & 1.000 & no \\
GPU server & MMLU & Phi-4 14B (bf16) & Qwen 3 30B-A3B (bf16) & 8/8 & 1.000 & 1.000 & no \\
\midrule
GPU server & GSM8K & Llama 3.3 70B (FP8) & Qwen 2.5 1.5B (bf16) & 41/4 & $<$0.001 & $<$0.001 & yes \\
GPU server & GSM8K & Llama 3.3 70B (FP8) & Llama 3.2 1B (bf16) & 36/4 & $<$0.001 & $<$0.001 & yes \\
GPU server & GSM8K & Llama 3.3 70B (FP8) & Phi-4 14B (bf16) & 21/0 & $<$0.001 & $<$0.001 & yes \\
GPU server & GSM8K & Llama 3.3 70B (FP8) & Ministral 3 8B (AWQ) & 27/2 & $<$0.001 & $<$0.001 & yes \\
GPU server & GSM8K & Llama 3.3 70B (FP8) & Ministral 3 3B (AWQ) & 29/4 & $<$0.001 & $<$0.001 & yes \\
GPU server & GSM8K & Llama 3.3 70B (FP8) & Qwen 3 4B (AWQ) & 25/3 & $<$0.001 & $<$0.001 & yes \\
GPU server & GSM8K & Llama 3.3 70B (FP8) & Qwen 2.5 7B (AWQ) & 26/5 & $<$0.001 & 0.002 & yes \\
GPU server & GSM8K & Llama 3.3 70B (FP8) & Qwen 3 30B-A3B (bf16) & 19/2 & $<$0.001 & 0.002 & yes \\
GPU server & GSM8K & Llama 3.3 70B (FP8) & Phi-4 Mini 3.8B (GPTQ-4) & 17/2 & $<$0.001 & 0.004 & yes \\
GPU server & GSM8K & Llama 3.3 70B (FP8) & Apertus 70B (bf16) & 18/3 & 0.001 & 0.007 & yes \\
GPU server & GSM8K & Llama 3.3 70B (FP8) & Ministral 3 14B & 23/6 & 0.002 & 0.009 & yes \\
GPU server & GSM8K & Llama 3.3 70B (FP8) & Phi-4 Mini 3.8B (bf16) & 13/4 & 0.049 & 0.147 & no \\
GPU server & GSM8K & Llama 3.3 70B (FP8) & Granite 4 H-Tiny (bf16) & 14/7 & 0.189 & 0.378 & no \\
GPU server & GSM8K & Llama 3.3 70B (FP8) & Kimi Linear 48B-A3B (bf16) & 12/7 & 0.359 & 0.378 & no \\
\midrule
GPU server & HellaSwag & Phi-4 14B (bf16) & Qwen 2.5 1.5B (bf16) & 62/2 & $<$0.001 & $<$0.001 & yes \\
GPU server & HellaSwag & Phi-4 14B (bf16) & Llama 3.2 1B (bf16) & 67/5 & $<$0.001 & $<$0.001 & yes \\
GPU server & HellaSwag & Phi-4 14B (bf16) & Kimi Linear 48B-A3B (bf16) & 66/5 & $<$0.001 & $<$0.001 & yes \\
GPU server & HellaSwag & Phi-4 14B (bf16) & Apertus 70B (bf16) & 50/6 & $<$0.001 & $<$0.001 & yes \\
GPU server & HellaSwag & Phi-4 14B (bf16) & Ministral 3 3B (AWQ) & 44/5 & $<$0.001 & $<$0.001 & yes \\
GPU server & HellaSwag & Phi-4 14B (bf16) & Phi-4 Mini 3.8B (bf16) & 37/5 & $<$0.001 & $<$0.001 & yes \\
GPU server & HellaSwag & Phi-4 14B (bf16) & Phi-4 Mini 3.8B (GPTQ-4) & 34/4 & $<$0.001 & $<$0.001 & yes \\
GPU server & HellaSwag & Phi-4 14B (bf16) & Qwen 2.5 7B (AWQ) & 30/4 & $<$0.001 & $<$0.001 & yes \\
GPU server & HellaSwag & Phi-4 14B (bf16) & Ministral 3 14B & 34/6 & $<$0.001 & $<$0.001 & yes \\
GPU server & HellaSwag & Phi-4 14B (bf16) & Granite 4 H-Tiny (bf16) & 27/6 & $<$0.001 & 0.002 & yes \\
GPU server & HellaSwag & Phi-4 14B (bf16) & Ministral 3 8B (AWQ) & 27/6 & $<$0.001 & 0.002 & yes \\
GPU server & HellaSwag & Phi-4 14B (bf16) & Qwen 3 4B (AWQ) & 22/6 & 0.004 & 0.011 & yes \\
GPU server & HellaSwag & Phi-4 14B (bf16) & Qwen 3 30B-A3B (bf16) & 12/6 & 0.238 & 0.476 & no \\
GPU server & HellaSwag & Phi-4 14B (bf16) & Llama 3.3 70B (FP8) & 8/7 & 1.000 & 1.000 & no \\
\end{longtable}

\endgroup

\end{document}
